\section{La phase d'accumulation}\label{sec:accumulation}
% Chapitre 5 (tache 10, passage chapitres 04-05). Figures: fig_gehalt,
% fig_jahre_sparquote, fig_kapital_zeit, fig_zinseszins, fig_sparbedarf,
% fig_start_verzoegerung (voir docs/figures.md). Aucun chiffre tape: macros de
% data/kennzahlen.tex (kapitel_5 et _macros_redaction_kapitel_5 dans
% scripts/rechnung_retraite.py). Controle croise (specification, section 8):
% l'age de base est \AlterYannBasis et l'annee \AnneeYannBasis (scenario de
% reference \ReferenzStrategie a \ReferenzSparquote %, revision du 30.09.2026),
% les memes macros que le resume; ce chapitre ne les recalcule jamais. Controle du
% salaire net: docs/salaire_controle.md (valeurs publiees reprises en constante
% SALAIRE_CONTROLE du script, avec leur source).

Ce chapitre r\'epond \`a la question~: \`a quel \^age ton \'epargne atteint-elle le capital du
chapitre~\ref{sec:capital}, selon la part de ton salaire que tu mets de c\^ot\'e et la strat\'egie que
tu choisis~? Le chapitre pr\'ec\'edent a fix\'e le capital \`a r\'eunir~; celui-ci le construit ann\'ee
apr\`es ann\'ee. La section~\ref{sec:accu-salaire} \'etablit ton salaire et le contr\^ole sur une table
publi\'ee, la section~\ref{sec:accu-taux} en tire ton \'epargne mensuelle, et la
section~\ref{sec:accu-age} donne l'\^age de d\'epart de chaque strat\'egie, y compris les combinaisons
qui n'aboutissent pas. Les sections suivantes mesurent l'effet du capital de d\'epart
(section~\ref{sec:accu-depart}), la part des int\'er\^ets compos\'es (section~\ref{sec:accu-interets}),
l'\'epargne qu'exigerait un d\'epart \`a 40, 45 ou 50~ans (section~\ref{sec:accu-besoin}) et le prix
d'un d\'emarrage plus tardif (section~\ref{sec:accu-retard}).

\subsection{Ton salaire d'ing\'enieur et son contr\^ole}\label{sec:accu-salaire}

Ton salaire r\'eel apr\`es le master n'est pas connu~; le plan le remplace par une r\'ef\'erence
publi\'ee et le suit sur toute ta carri\`ere.

\annahme{Salaire brut d'entr\'ee de \GehaltBrutEntree~€ par an, m\'ediane des ing\'enieurs de moins
d'un an d'exp\'erience\QH{einstiegsgehalt_brutto}, faute de grille Jungheinrich ou IG~Metall
exploitable. Progression r\'eelle de \GehaltWachstum~\% par an, d\'eduite de l'\'ecart entre
salaires d'entr\'ee et salaires apr\`es vingt-cinq ans d'exp\'erience dans la m\^eme
source\QH{gehaltssteigerung_real}~; c'est une coupe transversale, pas une carri\`ere observ\'ee, et le
plan la prolonge au-del\`a de ces vingt-cinq ans. Tout est en euros de 2026, avec le bar\`eme 2026
appliqu\'e chaque ann\'ee.}

\annahme{Salaire net approch\'e~: brut moins la part salariale des cotisations
sociales\QH{sv_arbeitnehmer}, moins l'imp\^ot sur le revenu du bar\`eme de 2026 sur un revenu
imposable r\'eduit du forfait de frais professionnels de \FraisProfessionnels~€, moins la
contribution de solidarit\'e~; classe d'imp\^ot~I, sans enfant, sans imp\^ot d'\'Eglise. Le mod\`ele
applique la contribution de solidarit\'e en entier d\`es que l'imp\^ot d\'epasse son seuil, sans la
zone de transition pr\'evue par la loi.}

La figure~\ref{fig:salaire-brut-net} montre la trajectoire qui en r\'esulte.

\linienfigur{fig_gehalt}[width=0.95\linewidth, xlabel={Ann\'ee}, ylabel={Salaire annuel (euros de 2026)},
  xticklabel style={/pgf/number format/set thousands separator={}}, scaled y ticks=false,
  yticklabel style={/pgf/number format/fixed}, legend pos=north west, legend cell align=left, legend style={fill=white}]%
  {jahr}{brutto,netto}%
  {{Brut},{Net}}%
  {Salaire annuel brut et net simul\'e, de la premi\`ere ann\'ee d'emploi \`a la fin de l'horizon du
  plan (euros de 2026).\label{fig:salaire-brut-net}}

Ce graphique montre un salaire net d'entr\'ee de \GehaltNetEntree~€ par an, soit
\GehaltNetEntreeMois~€ par mois et \GehaltNetPart~\% du brut, qui monte jusqu'\`a \GehaltNetFin~€ nets
pour \GehaltBrutFin~€ bruts en \AnneeHorizonFin. L'\'ecart entre les deux courbes s'\'elargit avec le
salaire, parce que l'imp\^ot est progressif. Le petit creux de la courbe nette en \GehaltCreuxAnnee{}
est un artefact du mod\`ele, d\^u \`a la contribution de solidarit\'e appliqu\'ee d'un coup~; il ne
change rien d'appr\'eciable au r\'esultat.

Ce calcul approch\'e a \'et\'e contr\^ol\'e une fois contre une table brut-net publi\'ee, sans retoucher
le mod\`ele pour la rejoindre\footnote{Deutschland-Rechner, \enquote{Brutto-Netto-Tabelle 2026}, classe
d'imp\^ot~I, assurance maladie publique au taux compl\'ementaire moyen, sans enfant, sans imp\^ot
d'\'Eglise, \href{https://www.deutschland-rechner.de/brutto-netto-tabelle}{deutschland-rechner.de},
consult\'e le 30~septembre 2026 (source secondaire).}. La table ne donne que des paliers ronds~; le
tableau~\ref{tab:controle-salaire-net} compare les deux paliers qui encadrent ton salaire d'entr\'ee.

\begin{table}[htbp]
  \centering\small
  \caption{Salaire net annuel publi\'e et calcul\'e par le mod\`ele, pour les deux paliers de la table
  qui encadrent le salaire d'entr\'ee (euros de 2026).}\label{tab:controle-salaire-net}
  \begin{tabular}{>{\bfseries}l r r r}
    \toprule
    \rowcolor[gray]{0.9}
    \textbf{Brut annuel} & \textbf{Net publi\'e} & \textbf{Net calcul\'e} & \textbf{\'Ecart relatif} \\
    \midrule
    \SalaireControleBrutBas~€ & \SalaireControlePublieBas~€ & \SalaireControleCalculeBas~€ & \SalaireControleEcartBas~\% \\
    \SalaireControleBrutHaut~€ & \SalaireControlePublieHaut~€ & \SalaireControleCalculeHaut~€ & \SalaireControleEcartHaut~\% \\
    \bottomrule
  \end{tabular}
\end{table}

Ce tableau montre un \'ecart de \SalaireControleEcartBas~\% et \SalaireControleEcartHaut~\%, bien sous
le seuil de \SalaireControleSeuil~\% fix\'e \`a l'avance. L'\'ecart va toujours dans le m\^eme sens~:
le mod\`ele est un peu plus g\'en\'ereux que la table, probablement parce que celle-ci suit la retenue
mensuelle sur salaire avec son propre forfait de pr\'evoyance, alors que le mod\`ele applique le bar\`eme
annuel \`a un revenu imposable simplifi\'e. Le d\'etail du contr\^ole est dans le fichier
\texttt{docs/salaire\_controle.md} du projet.

\subsection{Ton taux d'\'epargne}\label{sec:accu-taux}

\annahme{\'Epargne \`a partir de septembre 2027, \'egale \`a un pourcentage fixe du salaire net
(\SparquoteZehn, \SparquoteZwanzig, \SparquoteDreissig{} ou \SparquoteFuenfzig~\%), jamais moins de
\EpargneMinimum~€ par mois, et qui suit donc la hausse du salaire. Le mod\`ele compte la premi\`ere
ann\'ee enti\`ere, ce qui avance l\'eg\`erement le calendrier. Rendement r\'eel total de
\RenditeReal~\% par an pour les deux poches (chapitre~\ref{sec:depart}), constant chaque ann\'ee~;
dividendes de l'ETF maison impos\'es chaque ann\'ee et r\'einvestis nets, ETF monde capitalisant
impos\'e par la Vorabpauschale, pas de r\'e\'equilibrage entre les poches, ni frais de gestion ni frais
de courtage. \HypotheseBasculeReference. Le plafond annuel de la Vorabpauschale (Basisertrag) est
calcul\'e ici en euros r\'eels plut\^ot qu'en euros nominaux comme le pr\'evoit le droit, et elle n'est
plus mod\'elis\'ee apr\`es le d\'epart en retraite pour la r\'ef\'erence en retrait de \TauxRetrait~\%~:
les deux favorisent l\'eg\`erement l'ETF. Objectif atteint \`a la premi\`ere fin d'ann\'ee o\`u le revenu net de l'ann\'ee suivante
atteint \ZielMonat~€ par mois en euros de 2026.}

La premi\`ere ann\'ee, ces taux correspondent \`a \EpargneMoisZehn~€, \EpargneMoisZwanzig~€,
\EpargneMoisDreissig~€ et \EpargneMoisFuenfzig~€ par mois. D\`es le taux le plus bas, tu d\'epasses
donc le minimum de \EpargneMinimum~€ que tu t'es fix\'e, qui ne joue aucun r\^ole dans les calculs qui
suivent. Un rendement constant chaque ann\'ee est une simplification forte~: il ignore les krachs et
leur ordre, que le chapitre~\ref{sec:risques} r\'eintroduit par le rejeu historique et le tirage Monte Carlo.

\subsection{\^Age de d\'epart par strat\'egie et par taux}\label{sec:accu-age}

\annahme{Horizon de calcul de \HorizonJahre~ans, de 2027 \`a \AnneeHorizonFin, soit jusqu'\`a
\AlterHorizon~ans pour toi, avec le capital de d\'epart de \KapitalDepartMitte~€. \enquote{Non atteint}
signifie que l'objectif n'est pas atteint avant \AlterHorizon~ans. L'ETF distribuant vit de ses seuls
dividendes~; \HypotheseUmschichtung.}

Le tableau~\ref{tab:age-depart-strategies} croise les cinq strat\'egies du plan avec les quatre taux
d'\'epargne. C'est le r\'esultat central du document, et il est moins favorable que ce que l'on
esp\`ere en commen\c{c}ant un tel plan.

\begin{table}[htbp]
  \centering\small
  \caption{\^Age auquel l'objectif de \ZielMonat~€ nets par mois est atteint, selon la strat\'egie et
  le taux d'\'epargne (capital de d\'epart de \KapitalDepartMitte~€, euros de
  2026).}\label{tab:age-depart-strategies}
  \begin{tabularx}{\linewidth}{>{\bfseries}l *{5}{>{\centering\arraybackslash}X}}
    \toprule
    \rowcolor[gray]{0.9}
    \textbf{Taux d'\'epargne} & \textbf{Moiti\'e en maison} & \textbf{70~\% en maison}
      & \textbf{Tout en maison} & \textbf{ETF, retrait \TauxRetrait~\%} & \textbf{ETF distribuant} \\
    \midrule
    \SparquoteZehn~\% & \AlterMaisonFuenfzigZehn & \AlterMaisonSiebzigZehn & \AlterMaisonHundertZehn
      & \AlterEtfReferenzZehn & \AlterEtfUmschichtungZehn \\
    \SparquoteZwanzig~\% & \AlterMaisonFuenfzigZwanzig & \AlterMaisonSiebzigZwanzig
      & \AlterMaisonHundertZwanzig & \AlterEtfReferenzZwanzig & \AlterEtfUmschichtungZwanzig \\
    \SparquoteDreissig~\% & \AlterMaisonFuenfzigDreissig & \AlterMaisonSiebzigDreissig
      & \AlterMaisonHundertDreissig & \AlterEtfReferenzDreissig & \AlterEtfUmschichtungDreissig \\
    \SparquoteFuenfzig~\% & \AlterMaisonFuenfzigFuenfzig & \AlterMaisonSiebzigFuenfzig
      & \AlterMaisonHundertFuenfzig & \AlterEtfReferenzFuenfzig & \AlterEtfUmschichtungFuenfzig \\
    \bottomrule
  \end{tabularx}
\end{table}

Ce tableau se lit ligne par ligne. \`A \SparquoteZehn~\% d'\'epargne, aucune strat\'egie de dividendes
n'atteint l'objectif avant \AlterHorizon~ans~; seule la r\'ef\'erence y parvient, \`a
\AlterEtfReferenzZehn~ans. \`A \SparquoteZwanzig~\%, le constat reste le m\^eme~: \bk{aucune des trois
r\'epartitions en dividendes ni l'ETF distribuant n'atteint l'objectif dans l'horizon}. La
r\'epartition \ReferenzStrategie{} plafonne en \AnneeHorizonFin{} \`a \RevenuHorizonZwanzig~€ nets par
mois, \RevenuHorizonZwanzigPartCible~\% de la cible, alors que la r\'ef\'erence en ETF avec retrait de
\TauxRetrait~\% atteint l'objectif \`a \AlterEtfReferenzZwanzig~ans. C'est le r\'esultat le plus
d\'efavorable du document pour la th\`ese des dividendes, et c'est la raison pour laquelle le
sc\'enario de base \'epargne \ReferenzSparquote~\% et non \SparquoteZwanzig~\%.

\`A \ReferenzSparquote~\%, la r\'epartition \ReferenzStrategie{} atteint l'objectif \`a
\AlterYannBasis~ans, en \AnneeYannBasis, soit l'\^age l\'egal de la retraite, fix\'e \`a
\AgeLegal~ans\QH{regelaltersgrenze}~: \bk{au taux du sc\'enario de base, la strat\'egie des dividendes
ne t'offre aucune retraite anticip\'ee}. Au m\^eme taux, la r\'ef\'erence part \`a
\AlterEtfBasis~ans, \EcartAgeEtfMaisonBasis~ans plus t\^ot. C'est la traduction en \^age de l'\'ecart
de capital du chapitre~\ref{sec:capital}. Pour partir avant l'\^age l\'egal avec la r\'epartition
\ReferenzStrategie{}, il faut \'epargner au moins \QuoteAnticipeeMaisonSiebzig~\% de ton salaire net,
ce qui donne un d\'epart \`a \AlterAnticipeeMaisonSiebzig~ans sur la grille de la
figure~\ref{fig:annees-selon-epargne}, et \SparquoteFuenfzig~\% pour partir \`a
\AlterMaisonSiebzigFuenfzig~ans. Tout placer en maison ne fait gagner qu'une ann\'ee \`a
\SparquoteDreissig~\% (\AlterMaisonHundertDreissig~ans). L'autre voie est de changer de strat\'egie~:
la r\'ef\'erence en ETF part avant l'\^age l\'egal d\`es \QuoteAnticipeeEtfReferenz~\% d'\'epargne,
\`a \AlterAnticipeeEtfReferenz~ans.

La figure~\ref{fig:annees-selon-epargne} affine le tableau en faisant varier le taux d'\'epargne par
pas de cinq points~; une courbe absente \`a gauche signale un objectif non atteint dans l'horizon.

\linienfigur{fig_jahre_sparquote}[xlabel={Taux d'\'epargne (\% du salaire net)},
  ylabel={Ann\'ees d'\'epargne jusqu'\`a l'objectif},
  x filter/.expression={x*1e2}, unbounded coords=jump,
  table/empty cells with={nan}, legend pos=north east, legend cell align=left, legend style={fill=white}]%
  {quote}{maison50,maison70,\ColonneToutMaison,etf_ref}%
  {{Moiti\'e en maison},{70~\% en maison},{Tout en maison},{ETF, retrait \TauxRetrait~\%}}%
  {Nombre d'ann\'ees d'\'epargne, en comptant 2027, n\'ecessaires pour atteindre l'objectif selon le
  taux d'\'epargne, par strat\'egie (capital de d\'epart de \KapitalDepartMitte~€, horizon de
  \HorizonJahre~ans).\label{fig:annees-selon-epargne}}

Ce graphique montre le taux d'\'epargne minimal de chaque strat\'egie dans l'horizon~:
\QuoteMinEtfReferenz~\% pour la r\'ef\'erence, contre \QuoteMinMaisonFuenfzig{}, \QuoteMinMaisonSiebzig{}
et \QuoteMinMaisonHundert~\% pour la moiti\'e, 70~\% et la totalit\'e en maison. \`A ce taux minimal,
les strat\'egies de dividendes n'atteignent l'objectif qu'au bout de l'horizon ou presque, apr\`es
l'\^age l\'egal. Au-del\`a, chaque point d'\'epargne
suppl\'ementaire raccourcit de moins en moins la dur\'ee, parce que l'\'epargne vers\'ee compte de moins
en moins face aux gains du capital d\'ej\`a accumul\'e.

La figure~\ref{fig:capital-dans-temps} suit le capital de la r\'epartition de base ann\'ee par ann\'ee,
pour les quatre taux, face au capital requis du chapitre~\ref{sec:capital}.

\linienfigur{fig_kapital_zeit}[xlabel={Ann\'ee}, ylabel={Capital (millions d'euros de 2026)},
  xticklabel style={/pgf/number format/set thousands separator={}},
  y filter/.expression={y/1e6}, cycle list={{blue},{red},{brown},{black},{gray, dashed}},
  legend cell align=left, legend style={at={(0.5,-0.22)}, anchor=north, legend columns=3}]%
  {jahr}{q10,q20,q30,q50,ziel_kapital}%
  {{\'Epargne \SparquoteZehn~\%},{\'Epargne \SparquoteZwanzig~\%},{\'Epargne \SparquoteDreissig~\%},
   {\'Epargne \SparquoteFuenfzig~\%},{Capital requis du chapitre~\ref{sec:capital}}}%
  {Capital de la r\'epartition \ReferenzStrategie{} au fil des ann\'ees selon le taux d'\'epargne, et capital
  requis aux r\`egles de 2026 (euros de 2026).\label{fig:capital-dans-temps}}

Ce graphique montre qu'en \AnneeHorizonFin, le capital atteint \KapitalHorizonZehn~€ \`a
\SparquoteZehn~\% d'\'epargne, \KapitalHorizonZwanzig~€ \`a \SparquoteZwanzig~\%,
\KapitalHorizonDreissig~€ \`a \SparquoteDreissig~\% et \KapitalHorizonFuenfzig~€ \`a
\SparquoteFuenfzig~\%. Les courbes \`a \SparquoteZehn{} et \SparquoteZwanzig~\% ne croisent jamais la
ligne du capital requis, ce qui est le \enquote{non atteint} du tableau~\ref{tab:age-depart-strategies}.
La courbe du sc\'enario de base, \`a \ReferenzSparquote~\%, la franchit en \AnneeYannBasis{} avec
\KapitalYannBasis~€, soit \KapitalYannEcartCibleProzent~\% \KapitalYannSensCible{} capital du
chapitre~\ref{sec:capital}~: l'objectif se mesure en fin d'ann\'ee, et la derni\`ere ann\'ee d'\'epargne
fait passer le capital d'un peu moins \`a un peu plus que ce montant. Celle \`a \SparquoteFuenfzig~\%
la franchit d\`es \AnneeCroisementFuenfzig.

La simulation et le chapitre~\ref{sec:capital} visent le m\^eme capital parce qu'ils appliquent les
m\^emes seuils d'assurance maladie. Le plafond de l'assiette est revaloris\'e chaque ann\'ee avec les
salaires\QH{kv_schwellen_dynamisierung}~; le mod\`ele le tient donc constant en euros de 2026, et la
cotisation maladie et d\'ependance maximale vaut encore \KvPvMaxMoisObjectif~€ par mois en
\AnneeYannBasis, comme aujourd'hui. L'inflation n'all\`ege donc pas cette cotisation avec le temps.
\bk{Rien ne garantit toutefois que ce plafond suivra exactement les salaires.} Une r\'eforme qui le
rel\`everait alourdirait la cotisation et repousserait l'objectif~; \`a l'inverse, la
G\"unstigerpr\"ufung, que le plan n'int\`egre pas, all\'egerait l'imp\^ot. Les deux effets ne se
compensent pas n\'ecessairement~; retiens surtout qu'au taux du sc\'enario de base, le plan atteint
l'objectif \`a l'\^age l\'egal et n'a aucune marge pour partir plus t\^ot.

\subsection{Et avec \ZielNiedrigMonat~€ par mois~?}\label{sec:accu-ziel-niedrig}

\'Epargner la moiti\'e de ton salaire, c'est vivre avec l'autre moiti\'e~: une cible plus basse que
\ZielMonat~€ par mois est donc tout aussi pertinente \`a examiner. Le
tableau~\ref{tab:age-depart-ziel-niedrig} reprend exactement le tableau~\ref{tab:age-depart-strategies},
mais pour un objectif de \ZielNiedrigMonat~€ nets par mois.

\begin{table}[htbp]
  \centering\small
  \caption{\^Age auquel l'objectif de \ZielNiedrigMonat~€ nets par mois est atteint selon la
  strat\'egie et le taux d'\'epargne, et \'epargne mensuelle de la premi\`ere ann\'ee (capital de
  d\'epart de \KapitalDepartMitte~€, euros de 2026).}\label{tab:age-depart-ziel-niedrig}
  \begin{tabularx}{\linewidth}{>{\bfseries}l *{4}{>{\centering\arraybackslash}X}}
    \toprule
    \rowcolor[gray]{0.9}
    & \SparquoteZehn~\% & \SparquoteZwanzig~\% & \SparquoteDreissig~\% & \SparquoteFuenfzig~\% \\
    \midrule
    \tabellenkoerper{data/tab_ziel_niedrig}
    \bottomrule
  \end{tabularx}
\end{table}

La figure~\ref{fig:alter-ziel-vergleich} compare les deux cibles sur la m\^eme grille de taux
d'\'epargne, pour la r\'epartition \ReferenzStrategie{} et pour la r\'ef\'erence en ETF.

\linienfigur{fig_alter_ziel_vergleich}[xlabel={Taux d'\'epargne (\% du salaire net)},
  ylabel={\^Age d'atteinte de l'objectif (ans)},
  x filter/.expression={x*1e2}, unbounded coords=jump,
  table/empty cells with={nan}, legend pos=north east, legend cell align=left, legend style={fill=white}]%
  {quote}{maison70_ziel_haut,maison70_ziel_bas,etf_ziel_haut,etf_ziel_bas}%
  {{\ReferenzStrategie, \ZielMonat~€},{\ReferenzStrategie, \ZielNiedrigMonat~€},
   {ETF retrait \TauxRetrait~\%, \ZielMonat~€},{ETF retrait \TauxRetrait~\%, \ZielNiedrigMonat~€}}%
  {\^Age d'atteinte de l'objectif selon le taux d'\'epargne, pour \ZielMonat~€ et
  \ZielNiedrigMonat~€ nets par mois (capital de d\'epart de \KapitalDepartMitte~€, horizon de
  \HorizonJahre~ans).\label{fig:alter-ziel-vergleich}}

Baisser la cible \`a \ZielNiedrigMonat~€ rapproche le d\'epart de quelques ann\'ees \`a chaque taux
d'\'epargne, mais ne change rien \`a l'ordre des strat\'egies~: l'ETF avec retrait de
\TauxRetrait~\% reste le plus rapide, et il franchit l'\^age l\'egal \`a
\AlterEtfReferenzDreissigNiedrig~ans au taux de r\'ef\'erence de \ReferenzSparquote~\%, contre
\AlterMaisonSiebzigDreissigNiedrig~ans pour la r\'epartition \ReferenzStrategie{}. Le capital \`a
r\'eunir baisse dans les m\^emes proportions que la cible~: \KapitalNiedrigMaisonSiebzig~€ pour
\ReferenzStrategie{} et \KapitalNiedrigEtf~€ pour l'ETF seul, \`a comparer au capital du
chapitre~\ref{sec:capital} pour \ZielMonat~€ par mois.

\subsection{Le capital de d\'epart}\label{sec:accu-depart}

Ton capital de d\'epart se situe entre \KapitalDepartNiedrig~€ et \KapitalDepartHoch~€~; tous les
calculs pr\'ec\'edents retiennent le milieu, \KapitalDepartMitte~€. Pour le sc\'enario de base, ce
choix ne change pas l'\^age de d\'epart~: avec \KapitalDepartNiedrig~€, l'objectif est atteint \`a
\AlterBasisDepartNiedrig~ans, et avec \KapitalDepartHoch~€ \`a \AlterBasisDepartHoch~ans, comme avec
le milieu de la fourchette. L'\'ecart ne se voit qu'en fin d'horizon de calcul~: en \AnneeHorizonFin,
le capital atteint \KapitalHorizonDepartNiedrig~€ dans le premier cas et \KapitalHorizonDepartHoch~€
dans le second, soit \RevenuHorizonDepartNiedrig~€ et \RevenuHorizonDepartHoch~€ nets par mois.
Quelques milliers d'euros vers\'es en 2027 p\`esent donc peu face \`a trente ann\'ees et plus
d'\'epargne~; le taux d'\'epargne et la strat\'egie d\'ecident de l'issue, pas le capital de d\'epart.

\subsection{La part des int\'er\^ets compos\'es}\label{sec:accu-interets}

Dans le sc\'enario de base, l'essentiel du capital final ne vient pas de ton \'epargne. La
figure~\ref{fig:versements-contre-gains} s\'epare ce que tu as vers\'e de ce que le capital a
rapport\'e.

\linienfigur{fig_zinseszins}[xlabel={Ann\'ee}, ylabel={Montant cumul\'e (millions d'euros de 2026)},
  xticklabel style={/pgf/number format/set thousands separator={}},
  y filter/.expression={y/1e6}, legend pos=north west, legend cell align=left, legend style={fill=white}]%
  {jahr}{eingezahlt,gewinn}%
  {{Versements cumul\'es, capital de d\'epart compris},{Gains cumul\'es}}%
  {Versements cumul\'es et gains accumul\'es dans le sc\'enario de base (r\'epartition
  \ReferenzStrategie, \ReferenzSparquote~\% d'\'epargne, euros de 2026).\label{fig:versements-contre-gains}}

Ce graphique montre que les gains d\'epassent les versements \`a partir de \ZinsAnneeBascule. En
\AnneeYannBasis, sur un capital de \KapitalYannBasis~€, tu n'as vers\'e que \ZinsVersements~€~; les
\ZinsGains~€ restants, soit \ZinsPartGains~\% du capital, viennent des dividendes r\'einvestis et de la
hausse des cours, apr\`es imp\^ot. La courbe des gains s'acc\'el\`ere dans les derni\`eres ann\'ees,
quand le capital accumul\'e rapporte chaque ann\'ee plus que l'\'epargne vers\'ee. C'est ce qui rend
les premi\`eres ann\'ees d'\'epargne si pr\'ecieuses, et c'est aussi ce qui rend le r\'esultat si
sensible au rendement r\'eel suppos\'e~: la plus grande partie du capital d\'epend d'un rendement que
personne ne garantit.

\subsection{L'\'epargne n\'ecessaire pour partir \`a 40, 45 ou 50~ans}\label{sec:accu-besoin}

La question peut \^etre retourn\'ee~: combien faudrait-il \'epargner chaque mois pour partir \`a un
\^age choisi~? La figure~\ref{fig:epargne-selon-age} donne la r\'eponse pour la r\'epartition de base,
avec une \'epargne constante en euros de 2026.

\linienfigur{fig_sparbedarf}[xlabel={\^Age de d\'epart vis\'e (ans)},
  width=0.95\linewidth, ylabel={\'Epargne mensuelle (euros de 2026)}, ybar, bar width=28pt, xtick=data,
  enlarge x limits=0.3, ymin=0, enlarge y limits={upper, value=0.2}, scaled y ticks=false,
  nodes near coords, every node near coord/.append style={font=\scriptsize,
    /pgf/number format/fixed, /pgf/number format/precision=0}, legend pos=north east, legend cell align=left, legend style={fill=white}]%
  {alter}{monatlich}%
  {{\'Epargne mensuelle constante}}%
  {\'Epargne mensuelle constante n\'ecessaire pour atteindre l'objectif \`a l'\^age vis\'e,
  r\'epartition \ReferenzStrategie{} (capital de d\'epart de \KapitalDepartMitte~€, euros de
  2026).\label{fig:epargne-selon-age}}

Ce graphique montre qu'un d\'epart \`a 50~ans demande \SparbedarfCinquante~€ d'\'epargne par mois, soit
\SparbedarfCinquantePartNet~\% de ton salaire net d'entr\'ee. Pour 45~ans, il faudrait
\SparbedarfQuaranteCinq~€ par mois, \SparbedarfQuaranteCinqPartNet~\% de ce salaire, et pour 40~ans
\SparbedarfQuarante~€, \SparbedarfQuarantePartNet~\%. \bk{Aucun de ces trois \^ages n'est \`a port\'ee
d'un salaire d'ing\'enieur seul}~: m\^eme le d\'epart \`a 50~ans exige plus que la totalit\'e de ton
salaire net de d\'ebut de carri\`ere. Ils supposeraient un revenu nettement plus \'elev\'e, un
h\'eritage, ou une autre strat\'egie de revenu que les dividendes seuls.

\subsection{Commencer plus t\^ot ou plus tard}\label{sec:accu-retard}

La figure~\ref{fig:effet-depart-decale} d\'ecale d'un bloc le d\'ebut de ta vie active et de ton
\'epargne, toutes les autres hypoth\`eses restant celles du sc\'enario de base.

\linienfigur{fig_start_verzoegerung}[xlabel={D\'ecalage du d\'ebut de l'\'epargne (ann\'ees, n\'egatif
  si plus t\^ot)}, ylabel={\^Age d'atteinte de l'objectif (ans)}, xtick=data, legend pos=north west, legend cell align=left, legend style={fill=white}]%
  {verzoegerung_jahre}{ziel_alter}%
  {{\^Age d'atteinte}}%
  {\^Age auquel l'objectif est atteint quand l'\'epargne commence plus t\^ot ou plus tard que
  septembre 2027 (r\'epartition \ReferenzStrategie, \ReferenzSparquote~\% d'\'epargne, capital de
  d\'epart de \KapitalDepartMitte~€).\label{fig:effet-depart-decale}}

Ce graphique montre une droite de pente \VerzoegerungPente~: chaque ann\'ee de retard repousse ton
d\'epart d'une ann\'ee enti\`ere. Commencer \VerzoegerungAnsAvance~ans plus t\^ot aurait permis de
partir \`a \VerzoegerungAlterAvance~ans, commencer \VerzoegerungAnsRetard~ans plus tard repousse le
d\'epart \`a \VerzoegerungAlterRetard~ans. Le d\'ecalage vers le pass\'e n'est plus possible pour toi,
puisque ton master se termine en 2027~; il mesure ce que vaut une ann\'ee d'\'epargne. Le d\'ecalage
vers l'avenir, lui, est un risque r\'eel~: une ann\'ee de c\'esure, un emploi trouv\'e tard ou une
premi\`ere ann\'ee sans \'epargne co\^utent chacun une ann\'ee de retraite.

\bk{Retiens de ce chapitre qu'avec un salaire d'ing\'enieur, la strat\'egie des dividendes seuls
n'atteint jamais l'objectif \`a \SparquoteZwanzig~\% d'\'epargne, qu'\`a \ReferenzSparquote~\% elle ne
m\`ene qu'\`a l'\^age l\'egal, et qu'il faut environ \QuoteAnticipeeMaisonSiebzig~\% pour partir plus
t\^ot.} La r\'ef\'erence en ETF avec retrait de \TauxRetrait~\% part \EcartAgeEtfMaisonBasis~ans plus
t\^ot au taux du sc\'enario de base, et atteint l'objectif d\`es \SparquoteZwanzig~\%, \`a
\AlterEtfReferenzZwanzig~ans. Le chapitre~\ref{sec:portefeuille} examine l'ETF maison lui-m\^eme~: comment choisir ses
titres, et ce que change la r\'epartition entre ETF maison et ETF monde.

\Yannbox{\AnneeYannBasis}{En \'epargnant \ReferenzSparquote~\% de son salaire net, soit
\EpargneMoisBasis~€ par mois la premi\`ere ann\'ee, avec la r\'epartition \ReferenzStrategie{} et
\KapitalDepartMitte~€ de d\'epart, Yann atteint son objectif en \AnneeYannBasis, \`a
\AlterYannBasis~ans, l'\^age l\'egal de \AgeLegal~ans~: pas de retraite anticip\'ee. Sur les
\KapitalYannBasis~€ r\'eunis, il n'a vers\'e que \ZinsVersements~€~; les \ZinsGains~€ restants sont des
gains. Avec \KapitalDepartNiedrig~€ ou \KapitalDepartHoch~€ de d\'epart, l'\^age ne changerait pas.
\`A \SparquoteZwanzig~\% d'\'epargne, ses dividendes ne lui verseraient que \RevenuHorizonZwanzig~€
nets par mois en \AnneeHorizonFin. Pour partir plus t\^ot, il lui faudrait \'epargner
\QuoteAnticipeeMaisonSiebzig~\% (d\'epart \`a \AlterAnticipeeMaisonSiebzig~ans), ou passer \`a un ETF
monde avec un retrait de \TauxRetrait~\%~: \`a \AlterEtfBasis~ans au m\^eme taux d'\'epargne, et \`a
\AlterEtfReferenzZwanzig~ans d\'ej\`a avec \SparquoteZwanzig~\%.}
